\section{Invariant trace splittings and absorption}\label{sec:splittings}

We pass from Proposition~\ref{prop:uniform-unitaries} to exact
relations in \(A^\omega\). We first record the coordinate-selection
fact used in both passages.

\begin{lemma}[Coordinate selection]\label{lem:coordinate}
If \(x=[(x_j)]\in A^\omega\) has a bounded representative, then
\begin{align}
 \sup_{\tau\in\T_\omega(A)}\|x\|_{2,\tau}
   &=\lim_{j\to\omega}\sup_{\sigma\in\T(A)}\|x_j\|_{2,\sigma},
       \label{eq:coordinate-two}\\
 \sup_{\tau\in\T_\omega(A)}|\tau(x)|
   &=\lim_{j\to\omega}\sup_{\sigma\in\T(A)}|\sigma(x_j)|.
       \label{eq:coordinate-trace}
\end{align}
Consequently, finitely many such tests can be realized simultaneously
at a single coordinate with any prescribed positive slack above their
ultralimit bounds.
\end{lemma}

\begin{proof}
Every sequence of traces gives the upper bound in each identity.
Conversely, for each \(j\), choose \(\sigma_j\in\T(A)\)
within \(1/j\) of the supremum for the test in question.
Its limit trace realizes the reverse inequality; square roots and
absolute values commute with ultralimits of bounded scalar sequences.
For finitely many tests, the coordinates satisfying each individual
bound with positive slack form a set in \(\omega\). Their intersection
is again in \(\omega\), hence nonempty and indeed contains arbitrarily
large coordinates. The maximizing trace may be different for every test.
\end{proof}

\subsection{A weighted Haar unitary}

Keep the invariant separable algebra \(C\subset A^\omega\cap A'\)
from Proposition~\ref{prop:rows}. For \(a\in A_+\) and
\(\tau\in\T_\omega(A)\), define on \(C\)
\begin{equation}\label{eq:weighted-trace}
 \tau_a(c)=\tau(ac).
\end{equation}
This is a positive tracial functional. Positivity follows from
\(ac=a^{1/2}ca^{1/2}\) when \(c\geq0\), since \(a\) commutes with \(C\).
For \(c,d\in C\), traciality and the same commutation give
\(\tau(acd)=\tau(dac)=\tau(adc)\). Its norm is \(\tau(a)\).
If \(\tau(a)>0\), its normalization is therefore an element of
\(\T(C)\); if \(\tau(a)=0\), the functional is zero.
Proposition~\ref{prop:uniform-unitaries} thus gives, for its unitary \(V\),
\begin{equation}\label{eq:weighted-small-moments}
 |\tau(aV^\ell)|<\eps\,\tau(a)\leq\eps
 \quad
 \bigl(0\leq a\leq1,\ \tau\in\T_\omega(A),\
       1\leq\ell\leq m,\ \tau(a)>0\bigr),
\end{equation}
with zero left side when \(\tau(a)=0\).
The approximate invariance estimates apply to limit traces by
restriction to \(C\).

\begin{proposition}[A weighted Haar unitary]\label{prop:haar}
There is a unitary \(w\in(A^\omega\cap A')^{\bar\alpha}\) such that,
for every \(a\in A\), \(\tau\in\T_\omega(A)\), and
\(f\in C(\mathbb T)\),
\begin{equation}\label{eq:haar}
 \tau(af(w))=\tau(a)\int_{\mathbb T}f\,d\mu,
\end{equation}
where \(\mathbb T=\{z\in\C:|z|=1\}\) and \(\mu\) is normalized
Haar measure.
\end{proposition}

\begin{proof}
Choose increasing finite sets \(H_m\subset G\) whose union is \(G\),
and a sequence \(a_1,a_2,\ldots\) dense in the positive unit ball of
\(A\), with \(a_1=1\). Apply
Proposition~\ref{prop:uniform-unitaries} with
\((H,\eps,m)=(H_m,1/m,m)\), and denote its exponential unitary by
\(V^{(m)}\in C\).
It has a representative \((v^{(m)}_j)_j\) consisting of unitaries
in \(A\): lift its self-adjoint exponent to a bounded self-adjoint
sequence and exponentiate coordinatewise. The bounds for these
self-adjoint lifts may depend on \(m\), whereas the unitary
representatives always have norm one.

Since \(V^{(m)}\) commutes with \(A\), Lemma~\ref{lem:coordinate},
Equation~\eqref{eq:weighted-small-moments}, and the invariance
estimates allow a coordinate \(j(m)\geq m\) at which
\(w_m=v^{(m)}_{j(m)}\) satisfies
\begin{align}
 \|[w_m,a_r]\|_{2,u}&<2/m &&(1\leq r\leq m),\label{eq:first-diagonal}\\
 \|\alpha_h(w_m)-w_m\|_{2,u}&<2/m &&(h\in H_m),\notag\\
 \sup_{\sigma\in\T(A)}|\sigma(a_r w_m^\ell)|
   &<2/m &&(1\leq r,\ell\leq m).\notag
\end{align}
For the last line, apply Equation~\eqref{eq:coordinate-trace} to
the class \(a_r(V^{(m)})^\ell\). All coordinate tests are finite in
number, and their suprema already range over all traces.

The class \(w=[(w_m)]\) is a unitary. The first two lines of
Equation~\eqref{eq:first-diagonal}, density, and
\(\|[w_m,a-a_r]\|_{2,u}\leq2\|a-a_r\|\) show that it is central
relative to \(A\) and fixed by \(\bar\alpha\).
The last line, first for every \(a_r\), then by norm density for
positive contractions, and then by linearity for \(a\in A\), gives
\[
 \tau(aw^\ell)=0
 \qquad(a\in A,\ \tau\in\T_\omega(A),\ \ell\geq1).
\]
The same holds for negative powers, because
\(\tau(aw^{-\ell})=\overline{\tau(a^*w^\ell)}\).
The zeroth power has trace \(\tau(a)\).
Thus Equation~\eqref{eq:haar} holds for trigonometric polynomials.
Uniform approximation on \(\mathbb T\) and
\(|\tau(a f(w))|\leq\|a\|\|f\|_\infty\) prove it for every continuous
\(f\).
\end{proof}

\subsection{Equal trace projections}

\begin{proposition}[Equal trace projections]\label{prop:gamma}
The action \(\alpha\) has equivariant uniform property Gamma.
The projections in Definition~\ref{def:gamma} can be chosen to sum to \(1\).
\end{proposition}

\begin{proof}
Fix \(k\geq2\) and let \(w=[(w_j)]\) be the unitary of
Proposition~\ref{prop:haar}, represented by unitaries \(w_j\in A\).
Partition \(\mathbb T\) into \(k\) half-open arcs \(I_1,\ldots,I_k\)
of Haar measure \(1/k\), and write \(\chi_i=1_{I_i}\).
Choose continuous functions \(f_i^{(r)}:\mathbb T\to[0,1]\) satisfying
\[
 e_r:=\max_{1\leq i\leq k}\|f_i^{(r)}-\chi_i\|_{L^2(\mu)}
   \longrightarrow0.
\]
For example, replace each indicator near its two endpoints by linear
ramps, with the total lengths of the transition arcs tending to zero.
For this proof fix \(r\), and write \(f_i=f_i^{(r)}\) and \(e=e_r\).
Pointwise bounds and the triangle inequality imply
\begin{align}
 \|f_i^2-f_i\|_{L^2(\mu)}&\leq e,\label{eq:arc-bounds}\\
 \|f_i f_j\|_{L^2(\mu)}&\leq2e &&(i\ne j),\notag\\
 \Big\|\sum_i f_i-1\Big\|_{L^2(\mu)}&\leq ke,\notag\\
 \Big|\int_{\mathbb T} f_i\,d\mu-\frac1k\Big|&\leq e.\notag
\end{align}
For the first inequality, use
\(|f_i^2-f_i|\leq|f_i-\chi_i|\) almost everywhere.
For the second, expand
\(f_i f_j-\chi_i\chi_j=(f_i-\chi_i)f_j+\chi_i(f_j-\chi_j)\)
and use \(\chi_i\chi_j=0\) almost everywhere.
The remaining inequalities follow from \(\sum_i\chi_i=1\) almost
everywhere and the \(L^1\)--\(L^2\) inequality for a probability measure.

Set \(y_i^{(r)}=f_i^{(r)}(w)\).
These are central, invariant positive contractions, since continuous
functional calculus respects commutation and automorphisms.
Equation~\eqref{eq:haar} with \(a=1\), applied to the squared modulus
of each continuous function in Equation~\eqref{eq:arc-bounds},
gives the corresponding \(2\)-norm bounds for \(y_i^{(r)}\),
uniformly over \(\tau\in\T_\omega(A)\).
For every positive contraction \(a\in A\), it also gives
\begin{equation}\label{eq:arc-weighted}
 \left|\tau(a y_i^{(r)})-\frac{\tau(a)}k\right|
   =\tau(a)\left|\int_{\mathbb T}f_i^{(r)}\,d\mu-\frac1k\right|
   \leq e_r .
\end{equation}

Represent \(y_i^{(r)}\) by \(f_i^{(r)}(w_j)\).
Use \(H_r\) and \(a_1,\ldots,a_r\) from the proof of
Proposition~\ref{prop:haar}.
For this fixed \(r\), Lemma~\ref{lem:coordinate} gives a coordinate
\(j(r)\geq r\) such that the positive contractions
\(y_{i,r}=f_i^{(r)}(w_{j(r)})\) satisfy, simultaneously for all
indicated indices,
\begin{align}
 \|y_{i,r}^2-y_{i,r}\|_{2,u}&<e_r+1/r,\label{eq:second-diagonal}\\
 \|y_{i,r}y_{j,r}\|_{2,u}&<2e_r+1/r &&(i\ne j),\notag\\
 \Big\|\sum_i y_{i,r}-1\Big\|_{2,u}&<ke_r+1/r,\notag\\
 \|[y_{i,r},a_s]\|_{2,u}&<1/r &&(s\leq r),\notag\\
 \|\alpha_h(y_{i,r})-y_{i,r}\|_{2,u}&<1/r &&(h\in H_r),\notag\\
 \sup_{\sigma\in\T(A)}
 \left|\sigma(a_s y_{i,r})-\frac{\sigma(a_s)}k\right|
   &<e_r+1/r &&(s\leq r).\notag
\end{align}
This is a coordinate choice after \(f_i^{(r)}\) has been fixed.
It therefore needs no uniform continuity modulus, as \(r\) varies,
for these increasingly sharp arc approximations.

Let \(p_i=[(y_{i,r})_r]\in A^\omega\).
Their representatives are uniformly bounded positive contractions.
The first three lines of Equation~\eqref{eq:second-diagonal} imply
that \(p_i\) are orthogonal projections summing to \(1\); the next
two lines imply centrality and invariance by density.
For every sequence \((\sigma_r)\) in \(\T(A)\), the last line gives
\[
 \lim_{r\to\omega}\sigma_r(a_s y_{i,r})
   =\frac1k\lim_{r\to\omega}\sigma_r(a_s).
\]
This is Equation~\eqref{eq:gamma} for \(a_s\), for every limit trace.
Norm density and linearity extend it to all \(a\in A\).
Thus the projections satisfy Definition~\ref{def:gamma}.
\end{proof}

\subsection{The cocycle equations}

\begin{proof}[Proof of Theorem~\ref{thm:main}]
The trace hypotheses were checked in Section~\ref{sec:preliminaries},
and Proposition~\ref{prop:gamma} supplies equivariant uniform property
Gamma. All assumptions of Theorem~\ref{thm:absorption} hold.
It follows that \(\alpha\) and
\(\beta=\alpha\otimes\operatorname{id}_{\Zalg}\) are cocycle conjugate.
We spell out the orientation to obtain the exact equations of the
theorem.

An isomorphism \(\Phi:A\to B\) is automatically unital, since it is
surjective between unital algebras. Put
\(\gamma_g=\Phi\alpha_g\Phi^{-1}\).
If cocycle conjugacy is supplied in the form
\(\gamma_g=\Ad(u_g)\beta_g\) for a \(\beta\)-cocycle \(u\), then
Equations~\eqref{eq:cocycle} and~\eqref{eq:conjugacy} follow directly.
In the reverse orientation, suppose
\[
 \beta_g=\Ad(v_g)\gamma_g,\qquad
 v_{gh}=v_g\gamma_g(v_h).
\]
Here the unitaries lie in \(B\), whose multiplier algebra is \(B\)
itself. Set \(u_g=v_g^*\). Then
\[
 u_g\beta_g(u_h)
   =v_g^*v_g\gamma_g(v_h^*)v_g^*
   =\gamma_g(v_h^*)v_g^*
   =v_{gh}^*
   =u_{gh},
 \qquad
 \Ad(u_g)\beta_g=\gamma_g.
\]
Thus either orientation gives the required \(\beta\)-cocycle.
Its identity relation \(u_e=u_e^2\), together with unitarity,
forces \(u_e=1\). This proves both displayed equations and hence
the theorem.
\end{proof}
